% STATUS: reviewed
% ============================================================
%  [*] About this document — what it covers and how to read it
% ============================================================
\chapter*{Apie šį dokumentą}
\addcontentsline{toc}{chapter}{Apie šį dokumentą}
\markboth{Apie šį dokumentą}{Apie šį dokumentą}

Šiame dokumente aprašomas \product{} toks, koks jis buvo įdiegtas \docdate: ką daro kiekviena
jo dalis, kodėl ji sukurta būtent taip ir kaip dalys dera tarpusavyje. Dokumentas skirtas skaitytojui,
kuris išmano atitikties darbą ir finansų reguliavimą ir nori suprasti sistemą --- visų pirma
tai, kaip ji naudoja dirbtinį intelektą, --- neskaitydamas kodo. Pakeliui nurodomi failai ir
funkcijos, kad kodą atsivertęs skaitytojas rastų reikiamą vietą.

\section*{Kaip dokumentas sudarytas}

\begin{description}[style=nextline]
  \item[I dalis --- Sistemos apžvalga] Kam sistema skirta, kokios paslaugos ją sudaro ir
    kokius duomenis ji saugo (\cref{ch:intro,ch:arch,ch:data}).
  \item[II dalis --- Dirbtinis intelektas sistemoje] Ilgiausia dalis. Ji seka duomenis nuo
    teisės iki ataskaitos: kaip apskritai pasiekiamas modelis ir kiek kainuoja iškvietimas
    (\cref{ch:models}), kaip saugomos ir keičiamos instrukcijos (\cref{ch:prompts}), kaip
    teisės aktų tekstai tampa teisės aktų baze, kurioje galima ieškoti (\cref{ch:corpus}), ir
    pareigų registru (\cref{ch:obligations}), kas išskaitoma iš įmonės dokumentų
    (\cref{ch:company}), kaip rengiama atitikties ataskaita (\cref{ch:report}) ir kas su ja
    vyksta toliau (\cref{ch:register}), kaip rengiami atitikties programos dokumentų projektai
    (\cref{ch:programme}) ir kaip veikia pokalbių asistentas (\cref{ch:chat}). Po to
    \cref{ch:control} sutelkia tai, kas visa tai daro patikima: patikras, vertinimą,
    sprendžiančius žmones ir kainą.
  \item[III dalis --- Sistemos naudojimas ir eksploatavimas] Konsolė puslapis po puslapio
    (\cref{ch:console}), API (\cref{ch:api}), sauga (\cref{ch:security}), diegimas ir
    eksploatavimas (\cref{ch:ops}), testai (\cref{ch:testing}) ir tai, kaip sistema atliepia
    kursą, kuriam ji sukurta (\cref{ch:course}).
  \item[Žinynas] Visos instrukcijos (\cref{app:prompts}), visi pokalbio įrankiai
    (\cref{app:tools}), visi katalogo teisės aktai (\cref{app:acts}), parametrai
    (\cref{app:config}) ir terminų žodynėlis (\cref{app:glossary}).
\end{description}

\section*{Kas ką daro}

Kalbos modelius naudojančia sistema galima pasitikėti tik tiek, kiek aiški jos brėžiama riba
tarp to, ką parašė modelis, ką patikrino kodas ir ką nusprendė žmogus. Dokumentas šią ribą
visur žymi vienodai.

\begin{aibox}{Klausiama modelio}
Geltonas skydelis aprašo kalbos modelio iškvietimą: kas modeliui pateikiama, kas jam nurodoma
ir ką jis turi atsakyti. Paveiksluose geltonas stačiakampis yra modelio iškvietimas.
\end{aibox}

\begin{codebox}{Tikrina arba skaičiuoja kodas}
Mėlynas skydelis aprašo tai, kas atliekama be modelio: citata surandama jos dokumente,
nuostata randama teisės aktų bazėje, apskaičiuojamas balas, užtikrinamas taisyklės laikymasis.
Niekas tokiame skydelyje nepriklauso nuo to, ką modelis tąkart atsakė.
\end{codebox}

\begin{humanbox}{Sprendžia žmogus}
Rėmeliu apvestas skydelis žymi sprendimą, kurio sistema niekada nepriima pati: ataskaitos
apimties patvirtinimą, pareigos patvirtinimą, išvados būsenos pakeitimą, ataskaitos
patvirtinimą, parengto politikos projekto patvirtinimą.
\end{humanbox}

\begin{lessonbox}{Išmatuota}
Pilkame skydelyje pateikiamas matavimas arba nesėkmė, nurodant datą ir tai, kas dėl to buvo
pakeista. Šių skydelių skaičiai --- tai nurodytą dieną kurso pavyzdinės įmonės atveju atlikti
stebėjimai; tai nėra pažadai dėl kitų įmonių ar kitų dienų.
\end{lessonbox}

Kiekviena vieta, kurioje kviečiamas modelis, taip pat turi \emph{kelio kortelę}: lentelę su
tomis pačiomis dešimčia eilučių --- \emph{Paleidimas}, \emph{Įvestis}, \emph{Instrukcija ir
modelis}, \emph{Išvestis}, \emph{Kodo patikros}, \emph{Saugoma}, \emph{Nepavykus},
\emph{Kaina}, \emph{Žmogaus kontrolė} ir \emph{Kaip matuojama}. Jose nurodoma, kas kelią
paleidžia, ką modelis gauna, kokia instrukcija ir koks modelis naudojami, ką jis grąžina, ką
tikrina kodas, kas išsaugoma, kas nutinka nepavykus, kiek tai kainuoja, kur sprendžia žmogus ir
kaip kelias matuojamas. Kortelės išvardytos \cref{ch:control}.

\section*{Žymėjimas}

\begin{tabularx}{\textwidth}{@{}L{44mm}Y@{}}
  \toprule
  \hd{Užrašyta taip} & \hd{Reiškia} \\
  \midrule
  \file{analysis/pipeline.py} & Failas. API keliai nurodomi nuo \file{django/compliance/}, asistento paslaugos --- nuo \file{assistant/app/}, konsolės --- nuo \file{vite/app/src/}; visi kiti keliai nurodomi nuo repozitorijos šaknies. \\
  \code{run_prompt} & Pavadinimas kode: funkcija, konstanta, laukas, reikšmė. \\
  \prompt{report.gap_analysis} & Instrukcija pagal pavadinimą, kuriuo ji saugoma. \\
  \model{gpt-6.1-sol} & Modelis pagal pavadinimą, kurį jam suteikia DI tinklų sąsaja. \\
  \route{POST}{/api/reports} & API maršrutas. \\
  \ui{Veiksmai} & Užrašas taip, kaip jį rodo konsolė. \\
  instrukcijų injekcija (angl. \en{prompt injection}) & Pagrindinis DI terminas: pirmą kartą jį pavartojus skyriuje, nurodomas ir angliškas atitikmuo. \\
  9~sk., 9.3~posk., 3.2~pav., 2.1~lent. & Nuoroda į kitą dokumento vietą: pirmiausia numeris, po jo santrumpa; priedas nurodomas visu žodžiu, pvz., A~priedas. \\
  \status{met} \status{partial} \status{missing} \status{unknown} \status{na} & Penkios išvados būsenos: atitinka, iš dalies, neatitinka, neįrodyta, netaikoma. \\
  \bottomrule
\end{tabularx}

Sumos nurodomos JAV doleriais, pvz., 13,20~\$, nes taip sąnaudas skaičiuoja DI tinklų sąsaja. Laikas nurodomas pagal UTC. Pagrindinė produkto kalba --- lietuvių; jis veikia ir anglų kalba.
